\documentclass[UTF8, 11pt]{ctexart}

\usepackage[margin=1in]{geometry}
\usepackage{amsmath, amssymb}
\usepackage{booktabs}
\usepackage{hyperref}
\usepackage{url}
\usepackage{enumitem}
\setlist{leftmargin=*}
\setlength{\parskip}{0.35em}
\setlength{\parindent}{0pt}

\title{\textbf{阶段 3：复现与实现}\\
图划分算法实现与复现结果}
\author{小组成员：姓名 1，姓名 2，姓名 3}
\date{CS240：算法设计与分析，2026 年春季学期}

\begin{document}

\maketitle

\section{实现概览}

本项目实现为一个小型 Python 工程。代码位于 \texttt{src/} 目录下，按照图生成、
算法、评价指标、仿真、实验入口和绘图等功能拆分成多个模块。这样的模块化结构
便于独立测试每个组件，也方便在同一实验流水线中比较多种图划分方法。

主要文件包括：
\begin{itemize}
    \item \texttt{src/graphs.py}：构造 benchmark 图；
    \item \texttt{src/algorithms.py}：贪心、谱划分和 Kernighan--Lin 划分算法；
    \item \texttt{src/metrics.py}：划分质量评价指标；
    \item \texttt{src/simulation.py}：一致性平均仿真；
    \item \texttt{src/experiments.py}：实验运行与 CSV 导出；
    \item \texttt{src/plotting.py}：实验图像生成。
\end{itemize}

项目使用 Python、NumPy、NetworkX、pandas、matplotlib 和 pytest。所有报告中的
实验均使用固定随机种子，因此结果可以复现。

\section{实现的算法}

\subsection{贪心均衡划分}

贪心基线逐个将顶点分配到分区中。顶点按照加权度数从高到低处理，使影响较大的
顶点优先被放置。对于每个顶点，算法只考虑尚未达到容量上限的可行分区，并选择
相对于已分配邻居产生最小增量割值的分区。容量设为近似 $\lceil n/k\rceil$，
因此该方法可以保持接近完全均衡的分区大小。

该基线故意保持简单。它提供了快速比较点，有助于判断更复杂算法是否确实提升了
划分质量。

\subsection{递归谱二分}

谱方法反复二分当前最大的分区，直到得到 $k$ 个分区。对于一个子图，算法构造
加权邻接矩阵 $A$、度矩阵 $D$ 和图拉普拉斯矩阵 $L=D-A$。随后计算 $L$ 的特征
向量，并使用 Fiedler 向量对顶点排序。排序后的顶点列表在中点处切分。

该方法很适合作为较强的初始划分。在网格图和随机几何图中，Fiedler 向量通常
能够捕捉空间结构，因此谱二分往往能产生割值较低的平衡划分。

\subsection{Kernighan--Lin $k$ 路细化}

实现中的 Kernighan--Lin 方法从谱划分得到的初始解出发。随后执行若干轮全局
细化 pass。在每一轮中，算法遍历所有分区对。对于分区对 $(a,b)$，取出当前属于
这两个分区的顶点诱导子图，并以当前分配作为初始二分，调用 NetworkX 中稳定的
Kernighan--Lin 二分过程。

候选更新只有在不增加全局割边权重时才会被接受。该规则保证细化过程相对于割值
是单调不恶化的。如果完整一轮全局 pass 之后目标函数没有继续下降，算法提前
停止。

\section{复现范围}

原始 Kernighan--Lin 论文主要关注图二分。本项目复现的是其核心算法思想，而不是
历史论文中的全部实验设置。具体复现内容包括：
\begin{enumerate}
    \item 基于两个分区之间顶点交换的局部搜索细化；
    \item 重复执行 pass，直到没有明显割值改进；
    \item 将该方法作为细化过程，而不是仅作为随机初始划分器；
    \item 实验测量割值改进和额外运行时间。
\end{enumerate}

主要改写是 $k$ 路划分设置。我们没有只做一次全局二分，而是在 $k$ 路划分内部对
分区对执行成对 Kernighan--Lin 细化。这更适合分布式优化应用，因为目标系统包含
$k$ 个 工作节点，而不只是两个处理器。

\section{Benchmark 实例}

实现评估了四类图：
\begin{itemize}
    \item \textbf{网格图}：具有强局部几何结构的图；
    \item \textbf{Erdos--Renyi 随机图}：无明显结构的随机图；
    \item \textbf{随机几何图}：具有空间局部连接特性的图；
    \item \textbf{Zachary karate club 图}：小型真实社交网络。
\end{itemize}

benchmark 共包含八个实例：
\[
\begin{array}{lrrrr}
\toprule
\text{图类型} & \text{规模参数} & |V| & |E| & k \\
\midrule
\text{grid} & 6 & 36 & 60 & 4 \\
\text{grid} & 8 & 64 & 112 & 4 \\
\text{grid} & 10 & 100 & 180 & 4 \\
\text{erdos\_renyi} & 40 & 40 & 93 & 4 \\
\text{erdos\_renyi} & 80 & 80 & 248 & 4 \\
\text{geometric} & 50 & 50 & 330 & 4 \\
\text{geometric} & 90 & 90 & 628 & 4 \\
\text{karate} & 34 & 34 & 78 & 2 \\
\bottomrule
\end{array}
\]

\section{评价指标}

主要指标是割边权重：
\[
    \mathrm{cut}(p)=\sum_{(u,v)\in E,\;p(u)\neq p(v)}w_{uv}.
\]
割边权重越低，表示跨工作节点 依赖越少或越便宜。

第二个指标是归一化割，它将每个分区边界权重除以该分区体积后求和。这有助于
判断割值相对于各分区连接强度是否较小。

第三个指标是负载均衡比例：
\[
    \rho(p)=
    \frac{\max_i |\{v:p(v)=i\}|}{|V|/k}.
\]
越接近 1 表示工作负载越均衡。

第四个指标是运行时间，单位为秒，用于衡量每种划分算法的计算代价。

最后，一致性仿真报告最终一致性误差和跨分区消息数量。该仿真保持原始通信图
不变，因此不同划分不会改变收敛动力学；划分影响的是有多少消息需要跨工作节点
边界传输。

\section{复现结果}

实验复现了 Kernighan--Lin 风格局部细化的预期定性行为。在所有 benchmark 上，
平均割边权重为：
\[
    \text{greedy}=98.25,\qquad
    \text{spectral}=72.75,\qquad
    \text{Kernighan--Lin}=66.88.
\]
一致性仿真中的平均跨分区消息数为：
\[
    \text{greedy}=15460,\qquad
    \text{spectral}=11360,\qquad
    \text{Kernighan--Lin}=10440.
\]

Kernighan--Lin 细化方法取得了最低的平均割值和最低的平均跨工作节点 通信量。在
八个图实例中，它有六个实例取得最佳割值。这支持了预期复现目标：局部细化通常
可以相比简单基线提升划分质量，也能够在谱初始化基础上进一步改进或持平。

\section{偏差与失败原因分析}

与原论文相比，本项目有两个重要偏差。第一，原始 Kernighan--Lin 方法是二路图
划分算法，而本项目应用需要 $k$ 个分区。因此我们把成对二分细化作为组件嵌入到
$k$ 路划分框架中。

第二，本项目使用 NetworkX 中的成对 Kernighan--Lin 二分实现，而不是从零重写
每一个底层细节。这一选择提升了鲁棒性，使项目能够把重点放在算法比较、改写和
评估上。核心复现行为仍然是基于局部搜索的割值细化。

一致性实验也有一个限制：它测量跨工作节点 通信量，但没有改变底层通信拓扑。
因此，同一张图上不同划分的一致性误差几乎相同。这不是划分算法失败，而是仿真
设计导致的结果。该仿真隔离考察的是通信放置成本，而不是改变数值平均过程本身。

\section{运行方式}

可以用以下命令复现实验：
\[
\texttt{python src/experiments.py --output-dir outputs --seed 7}
\]
该命令会将原始 CSV 文件写入
\texttt{outputs/partition\_results.csv}，并将图像保存在
\texttt{outputs/figures/} 下。单元测试可以通过以下命令运行：
\[
\texttt{python -m pytest}
\]

\end{document}

